\documentclass[11pt]{article}
\usepackage[margin=1in]{geometry}
\usepackage{amsmath,amssymb,booktabs,graphicx,hyperref,xcolor,enumitem,caption}
\usepackage[numbers,sort]{natbib}
\hypersetup{colorlinks=true,linkcolor=blue!50!black,citecolor=blue!50!black,urlcolor=blue!50!black}
\newcommand{\pending}[1]{\textcolor{red!70!black}{[PENDING: #1]}}
\newcommand{\Cten}{C\textsubscript{1.0}}
\newcommand{\Cone}{C\textsubscript{0.1}}

\title{Birth of a Concept: Semantic Supervision, Weight Decay, and the\\
Emergence of Reusable Computation in Small Networks}
\author{Mnehmos research program\thanks{Experiments, analyses and drafting were carried out by an AI research agent
(Claude Code) under the owner's supervision; see the Provenance statement. Draft for owner review --- not submitted.}}
\date{Draft, \today}

\begin{document}
\maketitle

\begin{abstract}
We study when small networks (70K--130K parameters) learn a \emph{reusable} computation---one that applies to
input combinations never seen in training---rather than a lookup over trained combinations. In a registered
$4\times2$ factorial of 96 networks on a modular unit-conversion task, auxiliary supervision of the
\emph{canonical quantity} (the unit-invariant value) produced compositional generalisation to reserved
never-trained combinations only together with strong weight decay (reserved-test accuracy 0.655 vs.\ 0.101 for an
entropy-matched scrambled target; interaction $+0.48$, Holm $p=0.0012$). With weak decay the canonical
representation was equally decodable but not used compositionally. Using causal instruments validated on planted
mechanisms with known ground truth, we show that in the composing networks the canonical code is converted
correctly along (input, route) combinations that never occurred in training, and that the effect lives entirely
in the 12-dimensional subspace spanned by canonical-class means. In a second, historical study (36 networks with a
declared causal variable), training replays bit-exactly; networks with identical hyperparameters split into
near-perfect generalisers (OOD accuracy 0.998) and failures (0.26), and the split is all-or-none at the level of
the declared representation: the declared block either carries the canonical value of never-trained inputs
(interchange accuracy 0.99--1.00) or does not (0.09). \pending{mechanism and developmental results}
\end{abstract}

\section{Introduction}
Interpretability research often asks whether a network \emph{represents} a variable. For generalisation, the
question that matters is whether the network \emph{computes with} it on inputs it has not seen. We separate the two:
a representation can be decodable, even causally used on trained inputs, while the computation reading it remains
a lookup over trained combinations. We ask which training conditions make small networks acquire the reusable
computation, and what distinguishes networks that acquire it from otherwise identical networks that do not.

\paragraph{Contributions.}
\begin{enumerate}[leftmargin=*,itemsep=1pt]
\item A registered factorial (96 networks, 5,856 complete checkpoints every 100 steps, including optimizer state)
showing that canonical-quantity supervision promotes compositional generalisation only under strong weight decay.
\item A held-out causal instrument family (seven intervention conditions) validated on five planted mechanisms,
including the discriminating case of a correct canonical code read by a trained-combinations-only reader.
\item Exact (bitwise) replay of historical networks and their reconstruction with complete developmental records,
and an analysis of why identically configured networks generalise perfectly or fail.
\item Explicit limits: retrospective analyses are labelled exploratory; no hypothesis or threshold was changed
after seeing data; a validation-driven instrument change is reported as a deviation.
\end{enumerate}

\section{Tasks, networks and data}
\subsection{Modular unit conversion (factorial study)}
Quantities in two fictional dimensions (mass, acceleration) with three units each; conversion multiplies by unit
scale factors in $\mathbb{Z}_{13}^{\times}$: answer $= v\,s_u\,s_t^{-1} \bmod 13$. The canonical quantity is
$q=v\,s_u$. Twelve routes $(\text{dim},u,t)$ with $u\neq t$, 13 values: 156 combinations, of which 120 are trained,
18 development held-out and 18 reserved test (split hash \texttt{c103174ee0bfc1fa}); every withheld combination's
canonical quantity occurs in training via other routes. The student is a three-block transformer (129,280
inference parameters), AdamW lr $10^{-3}$, batch 128, 6,000 steps.

Populations: A narrow (42 combinations, one frame); B diverse; C diverse plus an auxiliary head predicting the
canonical quantity at the source-number token (after block 1); D the same head with an entropy-matched permuted
target. Weight decay 0.1 or 1.0; 12 paired seeds with identical initialisation across the eight conditions.

\subsection{Dimensional reasoning with a declared variable (historical study)}
Inputs \emph{value unit $\times$ value unit $\to$ force-unit $=$}; target $F = M A S \bmod 13$ with canonical
$M=m\,s_{mu}$, $A=a\,s_{au}$. Out-of-distribution (OOD) combinations: unit 2 with values 9--12 for either quantity;
their canonical values all occur in training. Groups: \texttt{mass\_only} and \texttt{acc\_only} train an
interchange-intervention objective \citep{geiger2022iit} declaring canonical $M$ (resp.\ $A$) in residual dims
0--15; \texttt{neither} has no such objective. 69.8K parameters, 6,000 steps, seeds 300--311 (doc 14/15).

\section{Measurement}
\subsection{Complete developmental records}
Every 100 steps: weights, AdamW moments and step counts, RNG states, and the exact decomposition of the summed
AdamW updates into the decoupled-decay part $-\eta\lambda\theta$ and the adaptive residual. Resumption is bitwise
exact (tested). The 96 factorial runs are preserved as per-run compressed archives with SHA-256 manifests
(15.5\,GB); owner-side verification of the delivered copies is pending.

\subsection{Held-out causal instruments}
Following interchange-intervention practice \citep{geiger2021causal,geiger2024causal}, we patch the residual at the
source-number token from a donor question (same dimension, different source unit) into a base question. A
\emph{number} code predicts the donor value converted along the base route; a \emph{canonical} code predicts the
donor quantity expressed in the base target unit, i.e.\ the answer for the base-route input $v' = q_d\,s_{u_b}^{-1}$.
Each pair is stratified by whether the base combination and the canonical-implied combination $(\text{route}_b,v')$
were trained or held out (reserved test combinations excluded everywhere). Following a pre-specified control that
failed on untrained networks, rates are reported relative to the no-intervention baseline (excess) and, after
planted-mechanism validation, as a \emph{switch rate}: among pairs whose unpatched prediction is not the target,
the fraction the patch moves onto it.

Seven conditions: I1 full residual after block 1; I2 number and unit tokens; I3 after block 2; I4 only the
subspace spanned by canonical-class means (fitted on trained bases); I5 its complement; I6 a same-canonical donor
(the canonical code predicts no change); I7 random subspaces of the same rank. Subspace patching can mislead
\citep{makelov2023illusion,wu2024reply}, and sufficiently flexible alignment maps can align networks with arbitrary
algorithms \citep{sutter2025nonlinear}; we therefore patch real donor activations, include I5/I7 controls, and
validate on planted mechanisms.

\paragraph{Planted-mechanism validation.} Five hand-built mechanisms with known ground truth
(Table~\ref{tab:planted}). All pass after one validation-driven change (switch rate instead of excess for P3: the
excess is bounded by $1-$baseline when an unpatched wrong answer coincides with the prediction). On real networks,
untrained (step-0) networks give $|\text{excess}|\le 0.02$ in every stratum and condition, and input-level token
swaps follow their known predictions.

\begin{table}[t]\centering\small
\caption{Planted-mechanism validation of the held-out causal instruments (all criteria fixed before running).}
\label{tab:planted}
\begin{tabular}{@{}llll@{}}\toprule
Planted mechanism & Ground truth & Criterion & Result\\\midrule
P1 canonical, compositional reader & converts $q$ on every route & I1 canonical $\ge 0.9$ all strata; I6 unchanged $\ge0.9$ & pass (1.00)\\
P2 number code, route lookup & number semantics & number $\ge 0.9$, canonical $\le 0.05$ & pass\\
P3 canonical code, trained-only reader & lookup over trained $(r,v')$ & switch $\ge0.9$ implied-trained; $\le0.10$ held-out & pass (1.00 / 0.00)\\
P4 rotated distributed code & as P1 & I4 $\ge0.9$; I5, I7 $\le 0.10$ & pass\\
P5 null & ignores the site & all $|\text{excess}|\le0.05$ & pass\\\bottomrule
\end{tabular}
\end{table}

\section{Registered factorial results}
\begin{table}[t]\centering\small
\caption{Reserved-test accuracy (72 items per network, best constant 0.167) and registered contrasts (paired by
seed, exact sign-flip, Holm over five). Independently re-derived from the archived weights with no discrepancy.}
\label{tab:primary}
\begin{tabular}{@{}lcccc@{}}\toprule
& A narrow & B diverse & C canonical aux & D scrambled aux\\\midrule
wd 0.1 & 0.061 & 0.098 & 0.134 & 0.058\\
wd 1.0 & 0.065 & 0.307 & \textbf{0.655} & 0.101\\\bottomrule
\end{tabular}\\[4pt]
\begin{tabular}{@{}lccc@{}}\toprule
Contrast & Mean & 95\% CI & $p_\text{Holm}$\\\midrule
H1 C$-$D @1.0 & $+0.554$ & [0.450, 0.659] & 0.0012\\
H2 interaction & $+0.478$ & [0.372, 0.587] & 0.0012\\
H3 C@1.0$-$C@0.1 & $+0.521$ & [0.431, 0.613] & 0.0012\\
H4 C$-$B @1.0 & $+0.348$ & [0.199, 0.479] & 0.0015\\
H5 B$-$A @1.0 & $+0.242$ & [0.176, 0.332] & 0.0012\\\bottomrule
\end{tabular}
\end{table}
All five registered contrasts were supported (Table~\ref{tab:primary}). Registered mechanistic predictions: M3
(downstream reads the code more canonically at strong decay) and M4 (decodable without composition at weak decay)
were confirmed; M1 (utilisation precedes composition) was directionally supported (9/3, $p=0.073$); M2's
seed-clustered re-analysis gives mean agreement 0.70 [0.60, 0.80], $p=0.003$.

\section{The canonical code is converted along never-trained combinations}
\pending{full 7-condition results across all 96 networks and developmental trajectories (I1 every 100 steps)}.
Pilot result on one network (\Cten, seed 2001, step 6,000): I1 switch rate 0.80 when the canonical-implied
combination was never trained (0.91 when trained); the 12-dimensional canonical-mean subspace alone (I4) reproduces
the full effect exactly, its complement (I5) gives 0.00, random 12-dimensional subspaces (I7) 0.02.

\section{Why semantic supervision needs strong weight decay}
\pending{H-A--H-E results (doc 41 \S41.4)}

\section{Historical networks: exact replay and all-or-none generalisation}
\subsection{Replay qualification}
The original training code (commit \texttt{aaa67c6}, exported unmodified) reproduced the historical networks of
seeds 304, 307, 310 and 311 \emph{bitwise}: every final parameter tensor equals the archived one and all 24 logged
loss/accuracy lines are byte-identical. \pending{instrumented replay (Q2), resume (Q3), reconstruction of all 36.}

\subsection{What separates generalisers from failures}
Among \texttt{mass\_only} networks that fit the training data (ID accuracy 1.0), OOD accuracy ranges from 0.25 to
0.998 under identical hyperparameters. A validated causal test separates them: interchanging the declared block
from a never-trained source into a trained base yields the counterfactual answer at 0.99/1.00 in the generalisers
(307, 311) and 0.09 in the failures (304, 310), and the declared variable's OOD accuracy is respectively 1.00 and
0.09 (chance 0.077). Probes fitted on trained inputs read 0.00 on OOD inputs even in the generalisers, a
false-negative illusion: the network represents unseen inputs in a different region that its readout still
decodes. \pending{developmental account, predictors}

\section{Limitations and scope}
Toy tasks (mod-13 arithmetic, fictional units), CPU-scale networks. Mechanistic analyses of the factorial are
retrospective and exploratory. The historical study was not designed for generalisation; its seed split is
observed, not manipulated; separating initialisation from data order would require a new registered experiment
(proposed in Appendix~\ref{app:prospective}, not run). Durable owner-side preservation of the checkpoint archives is
unverified at the time of writing.

\section*{Provenance}
All training, analysis and drafting were executed by an AI research agent (Claude Code) in a cloud container, under
the owner's registration, scope and approval rules. The specific model version serving each session is not
inferred from attribution metadata.

\appendix
\section{Proposed prospective experiment}\label{app:prospective}
\pending{Phase IV registration proposal}

\section{Reproducibility}
\pending{package description}

\bibliographystyle{plainnat}
\bibliography{refs}
\end{document}
